In this section, we describe the KG deployment infrastructure, including SPARQL query tools, constrained random walk, and optimizations.
We deploy Virtuoso 7.2.5 with Freebase. 
Entity types and names are specified via \texttt{type.object.type} and \texttt{type.object.name} with language tags. CVT nodes without readable names are flattened into two-hop relations (details are shown in section~\ref{sec:get_triples}).



\subsection{Entity Resolution}
\label{sec:entity_resolution}

Entity names are resolved to MIDs via exact match on \texttt{type.object.name@en}. Multiple candidates are disambiguated by selecting the entity with most types. A case-insensitive fallback is applied if needed. Results are cached. Listing~\ref{lst:entity_query} details this process.


\begin{lstlisting}[caption={Entity resolution query.},label=lst:entity_query,float=ht]
PREFIX ns: <http://rdf.freebase.com/ns/>
SELECT DISTINCT ?entity ?type WHERE {
  ?entity ns:type.object.name "Barack Obama"@en .
  ?entity ns:type.object.type ?type .
} LIMIT 100
\end{lstlisting}

\subsection{\texttt{get\_relations}}
\label{sec:get_relations}

\texttt{get\_relations(entity)} returns relation IDs. We split outgoing/incoming queries to avoid timeout, merge results, apply whitelist filtering, and optionally rank by BM25 similarity to return top-\textbf{k}. CVT flatten relations are appended. See Listing~\ref{lst:relations_query}.


\begin{lstlisting}[caption={Outgoing/incoming relation queries.},label=lst:relations_query,float=ht]
-- Outgoing
PREFIX ns: <http://rdf.freebase.com/ns/>
SELECT DISTINCT ?relation WHERE {
  ns:m.02mjmr ?relation ?tail .
  FILTER(isIRI(?relation) && STRSTARTS(STR(?relation), 
         "http://rdf.freebase.com/ns/"))
} LIMIT 100

-- Incoming
SELECT DISTINCT ?relation WHERE {
  ?head ?relation ns:m.02mjmr .
  FILTER(isIRI(?relation) && STRSTARTS(STR(?relation),
         "http://rdf.freebase.com/ns/"))
} LIMIT 100
\end{lstlisting}


\subsection{\texttt{get\_triples} and CVT Flattening}
\label{sec:get_triples}

\texttt{get\_triples(entity, relations)} returns triples. For each relation, we query tails/heads using \texttt{GROUP BY} and \texttt{SAMPLE} to obtain one name per node. CVT nodes (no name or ID-only name) are flattened: we collect their relations, form two-hop names (e.g., \texttt{rel1.rel2}), and resolve via two-hop queries. Flatten relations are cached for future \texttt{get\_relations} calls. 
Details are shown in Listings~\ref{lst:triples_query}--\ref{lst:cvt_query}.


\begin{lstlisting}[caption={Triples query with aggregation.},label=lst:triples_query,float=ht]
PREFIX ns: <http://rdf.freebase.com/ns/>
SELECT ?tail (SAMPLE(?name_en) AS ?preferred_name)
             (SAMPLE(?name_any) AS ?fallback_name)
WHERE {
  ns:m.02mjmr ns:people.person.place_of_birth ?tail .
  OPTIONAL { 
    ?tail ns:type.object.name ?name_en .
    FILTER(LANGMATCHES(LANG(?name_en), 'en'))
  }
  OPTIONAL { ?tail ns:type.object.name ?name_any . }
}
GROUP BY ?tail LIMIT 40
\end{lstlisting}

\begin{lstlisting}[caption={CVT relation query.},label=lst:cvt_query,float=ht]
PREFIX ns: <http://rdf.freebase.com/ns/>
SELECT DISTINCT ?relation WHERE {
  ns:<cvt_node_id> ?relation ?tail .
  FILTER(isIRI(?relation) && STRSTARTS(STR(?relation),
         "http://rdf.freebase.com/ns/"))
} LIMIT 50
\end{lstlisting}




